\booksection{Παράρτη$\mu$α: Ανάλυση επέκτασης $\mu$νή$\mu$ης σε Android κινητά}

\chapterintro{Όταν αγοράζετε κινητό, συχνά βλέπετε διαφη$\mu$ίσεις όπως «12GB+8GB RAM» — σαν να κερδίζετε 8GB RAM δωρεάν. Ό$\mu$ωςπράγ$\mu$ατι κάνει το κινητό πιο γρήγορο; Μπορεί αντίθετα να το βλάψει; Σε αυτή την ενότητα θα εξηγήσου$\mu$ε διεξοδικά την τεχνολογία «επέκταση $\mu$νή$\mu$ης».}

\booksubsection{I. Βασικός ορισ$\mu$ός}

Η “επέκταση $\mu$νή$\mu$ης” των Android κινητών (επίσης επέκταση $\mu$νή$\mu$ης, συγχώνευση $\mu$νή$\mu$ης, RAM Plus κ.λπ.)\textbf{Δεν προσθέτει πραγ$\mu$ατικά φυσικό τσιπ $\mu$νή$\mu$ης RAM}, αλλά $\mu$ια τεχνολογία λογισ$\mu$ικού σε επίπεδο συστή$\mu$ατος: δεσ$\mu$εύοντας δυνα$\mu$ικά $\mu$έρος του εσωτερικού αποθηκευτικού χώρου του τηλεφώνου (ROM/UFS flash) για να προσο$\mu$οιωθεί ως εικονική $\mu$νή$\mu$η RAM.

Με metaphor: η φυσική RAM είναι «γρήγορος χώρος εργασίας», ηεπέκταση $\mu$νή$\mu$η είναι «προσωρινή αποθήκη». Όταν ο χώρος εργασίας δεν επαρκεί, το σύστη$\mu$α $\mu$εταφέρει τα δεδο$\mu$ένα που δεν χρειάζονται στην αποθήκη και τα φέρνει πίσω όταν χρειάζονται. Η αποθήκη έχει $\mu$εγάλη χωρητικότητα, αλλά η $\mu$εταφορά είναι πιο αργή από το να τα έχετε στον χώρο εργασίας.

\booksubsection{II. Τεχνική αρχή}

Οι κύριοι κατασκευαστές υπάρχουν κυρίως δύο τρόποι χρήσης:

\subsubsection*{Συ$\mu$πίεση zRAM}

Δη$\mu$ιουργεί έναν χώρο συ$\mu$πιεσ$\mu$ένων δεδο$\mu$ένων $\mu$έσα στη RAM, συ$\mu$πιέζει τον κώδικα ανενεργών εφαρ$\mu$ογών $\mu$ε αναλογία περίπου 2:1. Τα δεδο$\mu$ένα παρα$\mu$ένουν στη RAM, ταχύτητα γρήγορη, δεν φθείρει το flash, αλλά η CPU συ$\mu$$\mu$ετέχει στη συ$\mu$πίεση/αποσυ$\mu$πίεση.

\subsubsection*{Δια$\mu$έρισ$\mu$α swap}

Δη$\mu$ιουργεί σταθερό χώρο στο UFS flash ως εικονική $\mu$νή$\mu$η, γράφει τα δεδο$\mu$ένα διεργασιών παρασκηνίου στην αποθήκευση. Εξαρτάται από τοεύρος ζώνης τουαποθήκευση, ταχύτητα πολύ χα$\mu$ηλότερη από τη φυσική $\mu$νή$\mu$η.

Η διαφή$\mu$ιση «12GB+8GB RAM» είναι συνήθως συνδυασ$\mu$ός 12GB φυσικής RAM + 8GB εικονικής $\mu$νή$\mu$ης flash. Ση$\mu$ειώστε ότι το zRAM προεπιλογή δεν καταλα$\mu$βάνει αποθηκευτικό χώρο, $\mu$όνο όταν το zRAM γε$\mu$ίσει εντελώς το σύστη$\mu$α δη$\mu$ιουργεί αρχείο swap στο storage, και αυτό σπάνια συ$\mu$βαίνει.

\booksubsection{III. Η διαφορά ταχύτητας είναι το βασικό bottleneck}

Ο βασικός περιορισ$\mu$ός της τεχνολογίας είναι η διαφορά ταχύτητας. Το ση$\mu$ερινό UFS 4.0 flash έχει διαδοχική ταχύτητα ανάγνωσης περίπου 4200MB/s, $\mu$όνο 1/10–1/15 της φυσικής $\mu$νή$\mu$ης LPDDR5X. Γι' αυτό στο ελαφρύ φορτίο του zRAM η ταχύτητα πλησιάζει τη φυσική RAM, αλλά σε υψηλό φορτίο η ταχύτητα πέφτει στα δεκάδες.

\booksubsection{IV. Σενάρια χρήσης και όρια οφέλους}

Ηεπέκταση $\mu$νή$\mu$η δενείναι πανάκεια, το αποτέλεσ$\mu$α διαφέρει ανά $\mu$οντέλο και σενάριο χρήσης:

\begin{itemize}
  \item \textbf{Χρήσι$\mu$ες περιπτώσεις}: συσκευές $\mu$ε 6GB φυσική RAM + UFS 3.1+ και $\geq$30GB ελεύθερο χώρο, $\mu$ε ενεργοποίηση 3GB επέκτασης, αυξάνουν τα ενεργά apps στο παρασκήνιο κατά 3-4, $\mu$ειώνοντας το ποσοστό επανεκκίνησης κατά 68\%, κατάλληλο για βαριά πολυδιεργασία (delivery, ταξί κ.ά.).
  \item \textbf{Καθη$\mu$ερινά σενάρια}: ση$\mu$αντική βελτίωση για πολλαπλές παρουσίες WeChat, διατήρηση καρτελών ιστού, απόκριση $\mu$ηνυ$\mu$άτων IM, ο χρόνος ψυχρής εκκίνησης εφαρ$\mu$ογών παρασκηνίου $\mu$ειώνεται αισθητά.
  \item \textbf{Περιορισ$\mu$ένα σενάρια}: περιορισ$\mu$ένη βελτίωση για σενάρια υψηλού εύρους ζώνης όπως φόρτωση $\mu$εγάλων παιχνιδιών, επεξεργασία βίντεο 4K.
  \item \textbf{Βλαβερές περιπτώσεις}: Σε $\mu$οντέλα $\mu$ε 8 GB RAM + UFS 2.2 και ανεπαρκή ελεύθερο χώρο, η ενεργοποίηση αυξάνει τις διακυ$\mu$άνσεις ρυθ$\mu$ού καρέ στα παιχνίδια, αυξάνει την κατανάλωση $\mu$παταρίας κατά την αναπαραγωγή σύντο$\mu$ων βίντεο και προκαλεί καθυστερήσεις στην εναλλαγή.
  \item \textbf{Άχρηστες περιπτώσεις}: για $\mu$οντέλα $\mu$ε πάνω από 16GB $\mu$νή$\mu$ης το όφελος είναι σχεδόν $\mu$ηδενικό, το πολύ να διατηρηθεί 1 επιπλέον εφαρ$\mu$ογή, αλλά αυξάνει το φορτίο CPU κατά 12\% και τη θερ$\mu$οκρασία κατά 2-3°C. Για $\mu$οντέλα $\mu$ε RAM πάνω από 12GB, συνιστάται απευθείας απενεργοποίηση για αποφυγή άσκοπων υπολογισ$\mu$ών συ$\mu$πίεσης.
\end{itemize}

\booksubsection{V. Ονο$\mu$ατοθεσία και τρόπος ενεργοποίησης ανά κατασκευαστή}

Διαφορετικές επωνυίες το λένε διαφορετικά: η Xiaomi «επέκταση $\mu$νή$\mu$ης», η Huawei «συγχώνευση $\mu$νή$\mu$ης», η Samsung «RAM Plus», η vivo «Extended RAM», η OnePlus «Memory Extension». Η λειτουργία ανήκει στις εγγενείς ικανότητες του συστή$\mu$ατος, δεν σχετίζεται ά$\mu$εσα $\mu$ε την έκδοση του Android.

Συνήθης διαδρο$\mu$ή ενεργοποίησης: [Ρυθ$\mu$ίσεις][Περισσότερες ρυθ$\mu$ίσεις]βρείτε το υπο$\mu$ενού [Μνή$\mu$η] ή [Απόδοση]ενεργοποιήστε την «Επέκταση $\mu$νή$\mu$ης» και επιλέξτε χωρητικότητα (συνήθως 2GB–7GB), $\mu$ετά από επανεκκίνηση ισχύει. Μετά το άνοιγ$\mu$α το σύστη$\mu$α ε$\mu$φανίζει συνδυασ$\mu$ό όπως «8GB+4GB».

\begin{tcbwarn}{Ση$\mu$ειώσεις}
\begin{itemize}
  \item Ηεπέκταση $\mu$νή$\mu$η καταλα$\mu$βάνει $\mu$έρος του αποθηκευτικού χώρου, εξασφαλίστε ελεύθερο χώρο τουλάχιστον 20-30GB, αλλιώς $\mu$πορεί να επηρεαστεί η ο$\mu$αλότητα του συστή$\mu$ατος.
  \item Το flash ως εικονική $\mu$νή$\mu$η έχει φθορά ανάγνωσης/εγγραφής, αν και τα σύγχρονα UFS flash έχουν $\mu$εγάλη διάρκεια ζωής, $\mu$ακροχρόνιο υψηλής συχνότητας swap $\mu$πορεί να επιταχύνει τη γήρανση τουαποθήκευση.
  \item Αν η φυσική $\mu$νή$\mu$η του κινητού είναι $\mu$ικρή (π.χ. 4GB και κάτω), ηεπέκταση $\mu$νή$\mu$η$\mu$πορεί $\mu$όνο ανακουφίζει προσωρινά, δεν λύνει ριζικά το πρόβλη$\mu$α του κολλή$\mu$ατος.
\end{itemize}

\end{tcbwarn}

\begin{tcbinfo}{Σύνοψη}
Ηεπέκταση $\mu$νή$\mu$ηείναι $\mu$ια ώρι$\mu$η τεχνολογία βελτιστοποίησης σε επίπεδο συστή$\mu$ατος, $\mu$πορείχωρίς αλλαγή υλικού να βελτιώσει τη multi-task ικανότητα και το ποσοστό διατήρηση ενεργού παρασκηνίου σε $\mu$εσαία-χα$\mu$ηλά $\mu$οντέλα, αλλά περιορίζεται από το τεράστιο χάσ$\mu$α ταχύτητας $\mu$εταξύ flash και φυσικής $\mu$νή$\mu$ης, δενείναι «πανακία». Αν θα την ανοίξετε και πόσο, πρέπει να κρίνετε βάσει του $\mu$εγέθους φυσικής $\mu$νή$\mu$ης, των προδιαγραφών flash, του ελεύθερου χώρου και των συνηθειών χρήσης.

\end{tcbinfo}

\begin{tcbpractice}{Πά$\mu$ε να δοκι$\mu$άσου$\mu$ε: Δείτε αν το κινητό σας έχει επέκταση $\mu$νή$\mu$ης}
\begin{enumerate}
  \item Ανοίξτε τις [Ρυθ$\mu$ίσεις], αναζητήστε «επέκταση $\mu$νή$\mu$ης», «συγχώνευση $\mu$νή$\mu$ης», «RAM Plus» ή «Extended RAM».
  \item Όταν τη βρείτε, ελέγξτε αν είναι ενεργοποιη$\mu$ένη και ποιαείναι η ρύθ$\mu$ιση χωρητικότητα επέκτασης.
  \item Συγκρίνετε $\mu$ε την ενότητα αυτή: πόσο $\mu$εγάληείναι η φυσική $\mu$νή$\mu$η του $\mu$οντέλου σας; Τι προδιαγραφές flash; Επαρκεί ο ελεύθερος χώρος;
  \item Δοκι$\mu$άστε να κλείσετε ή να $\mu$ειώσετε τηεπέκταση χωρητικότητας, παρατηρήστε για $\mu$ια εβδο$\mu$άδα τις αλλαγές στοδιατήρηση ενεργού παρασκηνίου και στη θερ$\mu$οκρασία, βρείτε τη πιο κατάλληλη ρύθ$\mu$ιση.
\end{enumerate}

\end{tcbpractice}
